\chapter{分治}

事实上分治仅仅是一个算法设计方法，并不特指某个算法。我们在这一章中介绍常见的想法和例子。

\section{标准序列分治}

均匀划分序列/点集来进行分治。

\begin{example}[可以用线段树维护的诸多信息]
    由于线段树是基于分治设计的数据结构，我们已经做了很多分治的练习了。
\end{example}

\begin{example}[平面最近点对]
    给定平面上的 $n$ 个点，查询距离最近的点对。
\end{example}
\begin{solution}
    先按 $x$ 排序。\texttt{solve(l,r)} 计算第 $l$ 到第 $r$ 个点之内的最近点对。首先递归 $\texttt{solve(l,m)}$ 和 $\texttt{solve(m+1,r)}$，得到两边的最近点对距离 $d$。接下来需要处理跨过中点的点对 $(i,j)$。注意到点对 $(i,j)$ 如果能更新答案，则有 $x_i\in [x_m-d,x_m],x_j\in[x_m,x_m+d]$，$y_j\in [y_i-d,y_i+d]$。也就是说，$j$ 被框在了 $[x_m,x_m+d]\times [y_i-d,y_i+d]$ 这个矩形内。另一方面，我们知道右侧任意两个点之间的距离都至少是 $d$，因此以每个点为圆心作半径为 $d/2-\epsilon$ 的圆，则所有这些圆之间都不交。因此，$[x_m,x_m+d]\times [y_i-d,y_i+d]$ 这个矩形内最多有 $O(1)$ 个点。

    那么做法就很明显了：按 $y$ 排序后，筛出 $x$ 在 $[x_m-d,x_m+d]$ 内的点，然后枚举左侧的每个点 $i$，逐个检查可能更新答案的 $j$。由于可能更新答案的点必须满足 $y_j\in [y_i-d,y_i+d]$，故这样的点在按 $y$ 排序并按 $x$ 筛选后位于序列的一段连续区间内，且长度为 $O(1)$。归并 $y$ 即可做到总复杂度 $O(n\log n)$。
\end{solution}

\begin{example}[CF1609F]
    查询有多少个区间的 $\max$ 和 $\min$ 的 \texttt{popcount} 相同。\texttt{popcount} 可以加强为任意函数 $f$。
\end{example}
\begin{solution}[做法一：序列分治]
    \texttt{solve(l,r)} 计算 $[l,r]$ 内的答案。首先递归 $\texttt{solve(l,m)}$ 和 $\texttt{solve(m+1,r)}$，得到两边的答案。接下来需要处理跨过中点的区间 $[i,j]$。令 $mx[i]=\max[i,m]$，$mx[j]=\max[m+1,j]$，$mi[i]=\min[i,m]$，$mi[j]=\min[m+1,j]$，则 $\max[i,j]=\max(mx[i],mx[j])$，$\min[i,j]=\min(mi[i],mi[j])$。那么问题就变成
    $$
    \sum_{i=1}^m\sum_{j=m+1}^n [f(\max(mx[i],mx[j]))=f(\min(mi[i],mi[j]))].
    $$
    接下来，枚举 $mx[i],mx[j]$ 以及 $mi[i],mi[j]$ 之间的大小关系，共有四种情况，再由对称性可以分为两类情况：
    \begin{itemize}
        \item 最值都在左侧：$mx[i]\ge mx[j],mi[i]\le mi[j]$。枚举 $i$，由于 $mx$ 与 $mi$ 都是单调的，满足这个限制的 $j$ 必然是一个前缀 $[m+1,j_i]$，且 $O(r-l+1)$ 扫一遍即可计算出所有 $j_i$。对于每个 $i$，满足这个限制的 $j$ 的数量为
        $$
        \sum_{j=m+1}^{j_i} [f(mx[i])=f(mi[i])]=(j_i-m)[f(mx[i])=f(mi[i])],
        $$
        $O(n)$ 计算即可。
        \item 最值都在右侧：$mx[i]<mx[j],mi[i]>mi[j]$ 同理。
        \item $\max$ 在左侧，$\min$ 在右侧：$mx[i]\ge mx[j],mi[i]>mi[j]$。枚举 $i$，满足这个限制的 $j$ 是一个区间 $[L_i,R_i]$，且 $O(r-l+1)$ 扫一遍即可计算出所有这些区间。对于每个 $i$，满足这个限制的 $j$ 的数量为
        $$
        \sum_{j=L_i}^{R_i} [f(mx[i])=f(mi[j])].
        $$
        可以注意到这是一个特殊的二维数点：对每个 $i$，查询满足 $j\in[L_i,R_i]$，$f(mi[j])=f(mx[i])$ 的 $j$ 的数量。差分后可以离线 $O(r-l+1)$ 完成（注意每个 $mi[j]$ 或者 $mx[j]$ 都是原序列中的某个 $a[i]$，因此在分治前离散化所有 $f(a[i])$ 即可将值域降到 $O(n)$）。
        \item $\max$ 在右侧，$\min$ 在左侧：$mx[i]<mx[j],mi[i]\le mi[j]$ 同理。
    \end{itemize}
\end{solution}
\begin{remark}
    这类题的通用做法都是枚举 $mx[i],mx[j]$ 与 $mi[i],mi[j]$ 之间的大小关系，然后变为偏序数点之类的问题。不过，某些题目可能具有特殊的性质，从而能够以更低的复杂度完成。
\end{remark}
\begin{solution}[做法二：扫描线]
    根据 $\max$ 的取值，所有的区间可以被表示为 $n$ 个不交矩形：对每个 $i$ 计算出 $i$ 能够作为区间最值的所有区间 $[l,r]$，这些区间可以由 $(\mathrm{pre}_i,i]\times [i,\mathrm{suf}_i)$ 刻画，其中 $\mathrm{pre}_i$ 和 $\mathrm{suf}_i$ 分别表示 $i$ 左/右侧第一个比 $a_i$ 大的元素的位置。同理，根据 $\min$ 的取值，也可以画出 $n$ 个不交矩形。

    枚举每一个值 $c$，考虑那些 $\max=c$ 的区间与 $\min=c$ 的区间对应的矩形。那么我们需要计算出\textbf{被覆盖了两次}的面积，扫描线即可做到 $O(\text{矩形个数}\log\text{矩形个数})$。由于总共的矩形数量是 $O(n)$，这样也是 $O(n\log n)$。
\end{solution}

\begin{exercise}
    统计有多少个区间满足 $\max-\min\le \mathrm{len}+k$
\end{exercise}

\begin{example}[P6109]
    先进行若干矩形加，然后查询若干次矩形 $\max$，点集为全平面。
\end{example}

\begin{example}[P3350]
    多询问网格图最短路
\end{example}

\section{非标准序列分治}

如果能保证分治的复杂度为 $T(n)=T(n_1)+T(n_2)+O(\min(n_1,n_2))$，那么复杂度也是 $T(n)=O(n\log n)$。这个结论可以通过将分治过程视为启发式合并的逆过程来证明。

\begin{example}[最值分治/笛卡尔树分治]
    P4755
\end{example}
\begin{solution}
    先预处理一个 RMQ，从而能够 $O(1)$ 回答任意区间的最值和最值位置。\texttt{solve(l,r)} 计算 $[l,r]$ 内的答案。按照 $[l,r]$ 之内的最值所在位置 $m$ 来分治，递归 $\texttt{solve(l,m-1)}$ 和 $\texttt{solve(m+1,r)}$ 来计算两边的答案。接下来需要处理跨过中点的区间 $[i,j]$。由于 $a_m$ 是 $[l,r]$ 内的最值，$\max [i,j]=a_m$。枚举 $[l,m]$ 和 $[m,r]$ 中较短的那一侧的每个 $i/j$，不妨设为 $[l,m]$，则需要询问 $[m,r]$ 中有多少个 $j$ 满足 $a_j\le a_m/a_i$。

    可以使用可持久化线段树来在线地回答这些询问，但是更好的方法是\textbf{二次离线}\graffiti{第一次离线是什么？}。首先差分询问变成 $[1,m-1]$ 内的和 $[1,r]$ 内的数点，我们用一个 $(l,m,t,c)$ 来表示：枚举所有 $i\in [l,m]$，查询 $[1,t]$ 内满足 $a_j\le a_m/a_i$ 的 $j$ 的数量，将结果乘 $c$ 后贡献到最终答案。这样的元组一共有 $O(n)$ 个，一共对应 $O(n\log n)$ 个询问。按照 $t$ 排序后就可以扫描线树状数组做到 $O(n\log^2 n)$ 甚至 $O(n\log^2 n/\log\log n)$，而空间是线性的。
\end{solution}
\begin{exercise}
    用前面标准序列分治的方法完成，需要做到同样的时空复杂度。
\end{exercise}
\begin{remark}
    实际上，这类问题能用最值分治做的似乎总能用标准序列分治做，而反过来则未必。不过最值分治可能好想好写一点。
\end{remark}

\begin{exercise}
    P9607, CF1156E
\end{exercise}

\begin{example}[CF475F]
    平面上给定 $n$ 个点，找到所有的极大点集满足：这些点集导出的矩形中，每一行和每一列都至少有一个点。
\end{example}

还有的时候，问题只会递归到一侧，或者递归到的子问题的总长度严格小于一个常数比例。此时复杂度通常是线性的。

\begin{example}[二分查找]
    二分查找可以视为是 $T(V)=T(V/2)+O(\text{判定复杂度})$ 的分治。
\end{example}
\begin{example}[快速选择算法]
    给 $n$ 个数，找出它们的第 $k$ 大。设计期望复杂度线性的算法。
\end{example}
\begin{exercise}[medians of medians]
    把序列每五个分成一组，取出每组的中位数，递归求出这些中位数的中位数作为划分点。证明这样做的复杂度是 $O(n)$。
\end{exercise}

\section{其他例子：多项式乘法和矩阵乘法}

我们再给出两个有趣的例子。

\begin{example}[Karatsuba 大整数乘法/多项式乘法]
    需要计算 $A(x)=\sum_{i=0}^{n-1} a_i x^i$ 和 $B(x)=\sum_{i=0}^{n-1} b_i x^i$ 的乘积。令 $X=x^{n/2}$，则问题变成计算 $(A_1(x)X+A_0)(B_1(x)X+B_0)$。直接递归需要计算四个乘积 $A_{0/1}B_{0/1}$，复杂度仍然是 $T(n)=4T(n/2)+O(n)=\Theta(n^2)$。但是注意到我们只需要计算出 $A_0B_0,A_1B_1$ 和 $(A_0+A_1)(B_0+B_1)$ 就可以表达出所需的全部四个乘积了。这样复杂度降低为 $T(n)=3T(n/2)+O(n)=\Theta(n^{\log_2 3})=O(n^{1.585})$。  
\end{example}
\begin{remark}
    这个神奇的构造可以认为是在\textbf{插值}：计算出 $(A_1X+A_0)(B_1X+B_0)$ 这个关于 $X$ 的二次多项式的任意三个点值就可以解出所有系数了。据此可以使用更多次的多项式来把复杂度降低至 $n2^{O(\sqrt{\log n})}$，这就是 Toom-Cook 算法。更多细节参见\href{https://www.cnblogs.com/Elegia/p/18020040/integer-multiplication}{EI 博客}。
\end{remark}

\begin{example}[Strassen 矩阵乘法]
    需要计算 $n$ 阶矩阵 $A,B$ 的乘积 $C$。将 $A,B$ 划分成四个 $n/2$ 阶的子矩阵，那么直接递归需要计算八个乘积 $A_{0/1}B_{0/1}$，复杂度仍然是 $T(n)=8T(n/2)+O(n^2)=\Theta(n^3)$。但是 Strassen 找到了构造：
    $$
    \begin{aligned}
        M_1 &= (A_{0,0} + A_{1,1}) \times (B_{0,0} + B_{1,1}) \\
        M_2 &= (A_{1,0} + A_{1,1}) \times B_{0,0} \\
        M_3 &= A_{0,0} \times (B_{0,1} - B_{1,1}) \\
        M_4 &= A_{1,1} \times (B_{1,0} - B_{0,0}) \\
        M_5 &= (A_{0,0} + A_{0,1}) \times B_{1,1} \\
        M_6 &= (A_{1,0} - A_{0,0}) \times (B_{0,0} + B_{0,1}) \\
        M_7 &= (A_{0,1} - A_{1,1}) \times (B_{1,0} + B_{1,1})
    \end{aligned}
    $$
    使得 $C$ 可以被这些矩阵表示：
    $$
        \begin{aligned}
        C_{0,0} &= M_1 + M_4 - M_5 + M_7 \\
        C_{0,1} &= M_3 + M_5 \\
        C_{1,0} &= M_2 + M_4 \\
        C_{1,1} &= M_1 - M_2 + M_3 + M_6
        \end{aligned}
    $$
    这样复杂度降低为 $T(n)=7T(n/2)+O(n^2)=\Theta(n^{\log_2 7})=O(n^{2.81})$。
\end{example}
\begin{remark}
    可以证明矩阵乘法的最优复杂度为 $n^{\omega+\epsilon}$，其中 $\omega=\inf\frac{\log r(\langle n,n,n\rangle)}{\log n}$ 是一个\textbf{常量}\graffiti{我猜 $\omega=2$。}。目前，人们知道 $\omega<2.371339$。矩阵乘法的复杂度是一个深刻的问题，参见 \href{https://en.wikipedia.org/wiki/Computational_complexity_of_matrix_multiplication}{wiki}。
\end{remark}

\section{偏序数点/CDQ 分治}

在前一章中我们提到过，可以用分治来查询偏序范围最值乃至任意分治信息。我们在此给出做法。

不妨假设范围是 $d$ 个正交半平面围出来的，也就是每个点是一个 $d$ 维带权点 $\mathbf x_i$，每次查询 $\mathbf x_i\le \mathbf y$ 的所有点的分治信息和\graffiti{由于这里没有顺序，需要分治信息可交换}。

我们对第一维进行分治：\texttt{solve(l,r)} 处理所有 $x_1,y_1\in[l,r]$ 的那些点与查询之间的贡献。那么 \texttt{solve(l,r)} 可以拆成三部分：\texttt{solve(l,m)}，\texttt{solve(m+1,r)}，以及处理所有 $x_1\in [l,m]$，$y_1\in [m+1,r]$ 的那些点与查询之间的贡献。注意到第三部分变成了一个相同规模但是维数减一的问题，可以递归到维数减一的问题。来算一下复杂度：设 $T(d,n,V)$ 是维数为 $d$，规模为 $n$，值域大小为 $V$ 时 \texttt{solve} 的复杂度，那么 $T(d,n,V)\le T(d,n_l,V/2)+T(d,n_r,V/2)+T(d-1,n,V_0)+O(n)$。

这里的初值是关键的。考虑 $d=1$ 的情况，如果直接按照分治的方法，那么我们需要 $O(n\log V)$ 的时间；而如果事先排好序了，那么就只需要 $O(n)$ 的时间了。所以，取决于点是否已经排好序了，$T(1,n,V)$ 为 $O(n)$ 或者 $O(n\log n)$。$d=2$ 的时候则是标准的二维数点，可以通过上一章中的扫描线方法做到 $O(n\log n)$。不过，我们在这里也介绍一种纯分治的方法。由于当点排好序时，$T(1,n,V)=O(n\log n)$，我们可以让 \texttt{solve(l,r)} 顺便支持将其内的所有点和询问按照第二维排序。这样，在递归处理完两个子问题之后，就可以 $O(n)$ 计算跨过中点的贡献，并且可以通过归并来 $O(n)$ 排序整个区间。这样我们也得到 $T(2,n,V)=T(2,n_l,V/2)+T(2,n_r,V/2)+O(n)=O(n\log V)$。对于更大的 $d$，使用前面的递归式即可算出 $T(d,n,V)=O(n\log^{d-1}V)$。另外，由于 $V$ 总是可以通过离散化变成 $O(n)$，我们可以在 $O(n\log^{d-1}n)$ 的时间内解决这个问题。

\begin{remark}
    当分治结构与询问无关时，将分治过程记录下来即可得到支持在线询问的数据结构，即树套树。所以，这个问题也存在 $O(\log^{d-1}n)$ 的在线做法，只不过这超出了大纲的范围。
\end{remark}

在实践中，$d$ 几乎总是不超过 $3$，更大的情况通常需要寻找其他方法\graffiti{比如数点的分块 bitset $O(n^2/w)$，一般情况的 KD tree $O(n^{2-1/d})$}。通常来说，分治方法在 $d=2$ 时的子问题的处理是最关键的。

\begin{exercise}[二维偏序]
    计算一个序列的逆序对数量。
\end{exercise}
\begin{exercise}[三维偏序]
    给定序列 $a,b,c$，查询满足 $a_i\le a_j,b_i\le b_j,c_i\le c_j$ 的 $(i,j)$ 的数量。
\end{exercise}
\begin{remark}
    如果只是求总对数的话，有一个容斥的 $O(n\log n)$ 做法。
\end{remark}
\begin{exercise}[带修二维数点]
    二维平面，支持插入或删除一个带权点 $\mathbf x$，查询所有 $\mathbf x\le \mathbf y$ 的点的分治信息和。
\end{exercise}
\begin{exercise}[全点集矩形加矩形求和]
    维护一个二维数组，支持矩形加，查询矩形和。
\end{exercise}
\begin{remark}
    如果不是全点集，那么这个问题是做不了的。实际上，行加列和就已经有矩阵乘法归约了。
\end{remark}
\begin{exercise}[P3157]
    每次删除一个数，查询剩下的序列的逆序对数
\end{exercise}
\begin{exercise}[P4169]
    二维平面，支持插入一个点，查询离一个点最近的点到它的距离，距离为曼哈顿距离 $|\Delta x|+|\Delta y|$。
\end{exercise}

\subsection{半在线}

我们的分治过程事实上可以支持一种比较弱的在线操作：假设点和询问初始全部已知，且在第一维上有序，但点权不是在最开始即给定的，而是在回答询问后才给出。比如说，给定 $(1,\mathbf x_1)$ 的点权，在回答询问 $(2,\mathbf y_2)$ 后再给定 $(3,\mathbf x_3)$ 的点权，然后再在回答询问 $(4,\mathbf y_4)$ 后给定 $(5,\mathbf x_5)$ 的点权，以此类推。由于计算跨过中点的贡献时只需要左边的点权全部已知，我们可以先递归左边，然后处理跨过中点的贡献，再递归右边。这在 OI 里通常被称为\textbf{cdq 分治}\graffiti{OI 中最早是陈丹琦推广的。但是没有人正式定义过 cdq 分治到底指的是什么。}。

\begin{exercise}[P2487]
    二维偏序意义下的 LDS，还需要计数并计数每个点在多少种方案中出现。
\end{exercise}

\begin{exercise}[P4093]
\end{exercise}

\begin{exercise}[在线卷积/分治 FFT]
    假设有一个支持 $O(n\log n)$ 计算多项式乘法的黑盒子，设计一个算法来在线地计算多项式乘法。也就是说，给定两个多项式 $A,B$，初始仅有 $A_0$ 和 $B_0$ 是已知的，计算出 $C_i=\sum_{j=0}^i A_j B_{i-j}$ 后才能获取 $A_{i+1}$ 和 $B_{i+1}$。
\end{exercise}
\begin{remark}
    实际上，由于我们可以 $O(n\log n)$ 计算 DFT，这个问题可以做得更好，参见\href{https://www.cnblogs.com/Elegia/p/vDH07-relaxed-multiplication.html}{EI 博客}。
\end{remark}

\section{整体二分}

通过批量地二分来复用预处理的结果，以降低复杂度。

\begin{example}[区间 kth]
    查询区间内第 $k$ 小的值。
\end{example}
\begin{solution}
    \texttt{solve(l,r,Q)} 实现：对于 $Q$ 中的每个询问 $[L,R,k,k_0]$，已知其结果在答案 $[l,r]$ 之间，且 $[L,R]$ 内小于 $l$ 的数有 $k_0$ 个，需要回答所有询问。
    
    我们对每个询问判断其答案是否 $\le m=(l+r)/2$，并求出新的 $k_0$，从而递归到子问题。那么只需要计算 $[L,R]$ 内有多少个数在值域 $[l,m]$ 中。将每个值在 $[l,m]$ 中的数视为 $1$，其他数视为 $0$，则问题变为一维数点。
    
    由于我们希望在 $\tilde O((r-l)+|Q|)$ 的复杂度内完成上面的过程，这个一维数点需要离散化或者使用数据结构。直接做带 $\log$，但是我们可以在递归的过程中顺便离散化，即同时传入当前值域对应的那些位置 $A$，并在递归过程中将询问的 $L,R$ 的规模降至 $|A|$。这样，复杂度就是 $T(V,|Q|,|A|)=T(V/2,|Q_l|,|A_l|)+T(V/2,|Q_r|,|A_r|)+O(|A|+|Q|+V)=O((|Q|+|A|)\log V)$ 了。
\end{solution}
\begin{remark}
    如果不做离散化，那么分治结构也是与询问无关的。此时，这种做法对应的在线数据结构是权值线段树套前缀和，想要做到 $O(\log n)$ 查询则需要分散层叠来避免多次二分定位 $L,R$。

    常见的在线做法则是前缀和套权值线段树，即在可持久化线段树的两个版本的差分上二分。
\end{remark}
\begin{exercise}
    查询区间内值域在 $[x,y]$ 中的第 $k$ 小值。
\end{exercise}
\begin{exercise}
    带修区间 $k$ 小。
\end{exercise}

\begin{example}[P4602]
    支持：插入删除一条带权树链，查询不包含某个点的树链的最大权值
\end{example}
\begin{solution}[做法一]
    维护每个点的答案，需要支持链补范围的插入删除，查询单点 $\max$。
\end{solution}
\begin{solution}[做法二]
    二分答案，需要判断是否存在一条链不包含该点。可以通过链加减、查单点来实现。
\end{solution}
\begin{solution}[做法三]
    注意到做法二中，只需要查询链交是否包含某个点。可以之间在线段树上二分。
\end{solution}

\begin{example}[P3527]
    环形序列 $a$，序列上每个位置有其所属的国家。支持 $a$ 区间加，对每个 $i$ 计算第 $i$ 个国家拥有的所有位置上的 $a$ 的和最早在何时达到 $p_i$。
\end{example}

\begin{example}[P5163]
   支持删除一条边，维护所有强连通分量。
\end{example}

\section{等价类分治}

处理一切范围修改查询问题\graffiti{传说中的 TB5。}。大体来说，如果这个问题放在序列区间上可以用标准的线段树完成，那么就可以扩大到任意\textbf{范围}上。问题可以被形式化表述如下：

\begin{definition}[范围修改查询]
    有 $n$ 个带权点 $a$ 和一个点的子集族 $\mathcal S$，支持：查询某个 $S\in \mathcal S$ 内的分治信息和 $f$，对某个 $S\in\mathcal S$ 内的点进行批量修改 $\Phi$，且能够进行快速的
    \begin{itemize}
        \item[信息合并] 可以根据 $f(a;S_1)$ 和 $f(a;S_2)$ 快速计算 $f(a;S_1\uplus S_2)$
        \item[修改合并] 可以快速计算 $\Phi_1\circ \Phi_2$
        \item[修改分裂] $S=S_1\uplus S_2$ 上的修改 $\Phi(a;S)$ 可以快速分解为子集 $S_1$ 和 $S_2$ 上的修改 $\Phi(a;S_1)$ 和 $\Phi(a;S_2)$ 
        \item[应用修改] 可以根据 $f(a;S)$ 快速计算 $f(\Phi(a;S);S)$  
    \end{itemize}
\end{definition}

我们首先对于给定的修改与查询序列定义点的等价关系与等价类。

\begin{definition}[等价关系与等价类]
    如果两个点 $u,v$ 满足：涉及到 $u$ 的修改/查询与涉及到 $v$ 的修改/查询完全相同，那么我们称 $u,v$ 是等价的。
\end{definition}

显然，同一个等价类内的点可以缩到一起。这样，一次修改与查询的复杂度就只与等价类的数量有关。而修改与查询的数量越少，等价类的数量也就越少。这样，我们可以设计一个关于时间的分治过程。记 $\mathcal C[l,r]$ 为 $[l,r]$ 时间内的修改与查询导出的所有等价类。如果我们只需要做 $[l,r]$ 时间内的修改与查询，那么我们就只需要关注 $\mathcal C[l,r]$ 中的每个等价类的信息和修改情况。

对于每个等价类 $C_i\in \mathcal C[l,r]$，我们希望在 $\mathtt{solve}(l,r)$ 开始前就知道其总信息 $f(a;C_i)$，并且在函数结束知道 $[l,r]$ 这段时间内的修改在其上的总修改 $\Phi(a;C_i)$。自然，我们需要利用 $\mathtt{solve}(l,m)$ 和 $\mathtt{solve}(m+1,r)$ 来完成。由于子区间的等价类集合相当于 $\mathcal C[l,r]$ 中的等价类进行一些合并得到的，当我们想要递归进某一边的之前，就可以对 $\mathcal C[l,r]$ 内的等价类做一些合并来得到 $\mathcal C[l,m]$ 或者 $\mathcal C[m+1,r]$，并且做相应的信息合并。当递归完 $[l,m]$ 或者 $[m+1,r]$ 回到 $[l,r]$ 时，我们需要重新分裂先前合并的等价类，并将 $\mathcal C[l,m]$ 或者 $\mathcal C[m+1,r]$ 中每个等价类上的总修改标记分裂到 $\mathcal C[l,r]$ 中的对应等价类上。 

\begin{descbox}{算法}[等价类分治]
    $\mathtt{solve}(l,r,\mathcal C[l,r])$ 实现以下功能：在函数开始前，需要保证 $\mathcal C[l,r]$ 即为 $[l,r]$ 中的修改查询导出的所有等价类，每个等价类 $C_i$ 由其中的任一代表元 $\alpha_i$ 表示，且其总信息 $f(a;C_i)$ 已经计算完毕。在函数结束后，计算出 $[l,r]$ 内的所有修改在每个等价类 $C_i$ 上的总修改量 $\Phi(a;C_i)$，以及修改后得到的新的总信息 $f(\Phi(a;C_i);C_i)$。
    
    $\mathtt{solve}(l,r,\mathcal C[l,r])$ 的具体实现如下。令 $m=(l+r)/2$。首先，通过\textbf{某种方法}计算出 $\mathcal C[l,m]$，以及每个 $C_i\in \mathcal C[l,r]$ 在 $\mathcal C[l,m]$ 中位于哪个等价类。据此可以利用每个 $C_i\in\mathcal C[l,r]$ 的总信息 $f(a;C_i)$ 来合并计算出每个 $C_j\in\mathcal C[l,m]$ 的总信息 $f(a;C_j)$。接下来，递归 $\mathtt{solve}(l,m,\mathcal C[l,m])$ 计算出 $[l,m]$ 内的所有修改在每个 $C_j\in\mathcal C[l,m]$ 上的总修改量 $\Phi(a;C_j)$。然后，将每个 $\Phi(a;C_j)$ 分解到 $\mathcal C[l,r]$ 中对应的等价类上，从而得到每个 $C_i\in \mathcal C[l,r]$ 上的总修改 $\Phi(a;C_i)$。将每个 $C_i\in\mathcal C[l,r]$ 上的总修改 $\Phi(a;C_i)$ 应用到对应的信息 $f(a;C_i)$ 上即可得到 $[l,m]$ 内的所有修改执行完毕后，每个等价类 $C_i$ 中新的总信息 $f(\Phi(a;C_i);C_i)$。然后同样地对 $[m+1,r]$ 执行以上操作。
\end{descbox}

假设 $i$ 个范围至多导出 $P(i)$ 个等价类，算法中用于计算等价类划分的\textbf{某种方法}面对 $n$ 个点和 $m$ 个范围时具有 $D(m,n)$ 的复杂度，则分治的复杂度为
$$
T(m)=2T(m/2)+O(P(m))+D(m/2,P(m)).
$$
而在 $\mathtt{solve}(1,m,\mathcal C[1,m])$ 开始前，我们还需要花费 $D(m,n)$ 的时间来计算 $\mathcal C[1,m]$，因此总时间为 $T(m)+D(m,n)$。

注意我们还可以对操作序列按 $B$ 大小分块，这样总共就是 $\frac{m}{B}(T(B)+D(B,n))$ 的复杂度。

\begin{remark}
    实践上，使用等价类分治时最困难的地方基本在于\textbf{某种方法}。如果忽略这个部分的话，可以证明这个算法的复杂度是理论最优的 $\Theta(m\cdot \max_{1\le i\le m}P(i)/i)$。
\end{remark}

\begin{example}[常规线段树问题]
    $m$ 个区间至多划分出 $\min(n,2m+1)$ 个等价类，且按端点排序后计算这些等价类只需要 $O(n+m)$ 的时间。所以在最外部排一次序后就有 $T(m)=O(m\log \min(n,m))$，$D(m,n)=O(n+m)$，总复杂度 $O(m\log \min(n,m)+n)$（而不是 $n+m\log n$！）
\end{example}

\begin{example}[矩形加矩形和]
    $m$ 个矩形至多划分出 $O(\min(m^2,n))$ 个等价类。在外层将询问按照 $\sqrt{n}$ 分块，则可以假设 $m\le \sqrt{n}$。此时，在最外部将两维分别排序后，可以在线性时间内计算所有等价类（行等价类与列等价类的笛卡尔积）。所以 $D(m,n)=O(m^2+n)$，$T(m)=O(m^2)$，总共 $O(n+q\sqrt{n})$。
\end{example}
\begin{remark}
    半平面加半平面求和、圆加圆求和也是相同的复杂度（忽略 $\log$ 因子）。
\end{remark}

\begin{example}[树链加树链和]
    $m$ 个树链可以按照虚树的方式划分出 $O(\min(n,m))$ 个等价类，所以是 $O(m\log\min(n,m)+n)$ 的。
\end{example}

在有位移操作（如区间翻转和平移）时，等价类分治仍然可以应用。只不过此时运行完一侧的修改后，除了每个等价类内部的总修改以外，还可能会发生等价类之间的重新排列。此时，可以先按照函数复合的方式自下而上地计算出等价类的划分，再运行标准的等价类分治。

\begin{example}[文艺平衡树]
    
\end{example}

\begin{example}[*P7723]
    TB5x
\end{example}